\textbf{Object of the study}: the process of automating personal task management using artificial intelligence.

\textbf{Subject of the study}: design and implementation of an AI assistant that autonomously organizes user tasks using advanced time management methodologies and presents them in an intuitive form, removing the need for the user to delve into their complexity.

\textbf{Goal of the work}: design and implementation of a minimum viable product (MVP) of a personal AI assistant that intelligently manages user tasks using a comprehensive methodology (in particular, GTD) “under the hood”, thereby radically simplifying the process of planning and achieving goals for the user.

\textbf{Research and development methods}: analysis of existing task management systems and productivity methodologies (including GTD); systems design; object-oriented programming; development of microservice architecture; use of the C++ (with the userver framework) and Python (with the python-telegram-bot library) programming languages; application of artificial intelligence technologies (YandexGPT); database management (PostgreSQL); containerization (Docker, Docker Compose); software testing (modular, integration, user); deployment in cloud infrastructure (Yandex Cloud).

\textbf{Results and their novelty}:
During the work, the problem of the lack of an intuitive AI solution for task management that effectively implements a complex methodology (GTD), completely hiding its complexity from the user, was solved.

\begin{enumerate}
\item backend services in C++/userver, implementing the main GTD lists (current-actions, sometime-later, waiting, notes, calendar);

\item frontend service in the form of a Telegram bot in Python;

\item ai-proxy service in Python for integration with YandexGPT.

\end{enumerate}

The system provides automatic processing and AI categorization of tasks, task management, report generation and event export in .ics format, while the complexity of the GTD methodology is completely abstracted from the user. The scientific novelty lies in the proposed and implemented concept of an AI assistant that autonomously implements an advanced time management methodology (GTD) «under the hood», reducing the entry threshold for its use.

\textbf{Practical significance}:

\begin{enumerate}
\item An MVP of an AI assistant has been developed that can automate routine task management operations and help users cope with cognitive overload.

\item A flexible and scalable technological base (microservices, C++, Python, YandexGPT, PostgreSQL, Docker, Yandex Cloud) has been created for further development of the product.

\item An approach to simplifying the use of complex productivity methodologies through AI is proposed and tested.

\item The materials of the VKRB can be used for educational purposes in studying the design and development of AI systems and microservice applications.
\end{enumerate}

\textbf{Application area}: personal use to improve productivity and organize tasks; potential for integration into corporate systems.